\documentclass[12pt, a4paper, oneside]{ctexart}
\usepackage{amsmath,amsthm,amssymb,graphicx}
\usepackage{float}
\usepackage[margin=2.5cm]{geometry}
\usepackage{xcolor}
\usepackage{enumitem}
\begin{document}
\section*{子域和扩域}

定义：给定域$F$，若$F$的子集$K$在$F$的运算下也构成一个域，则我们称$K$为$F$的子域，也称$F$为$K$的扩域。

不难发现，域$F$是其本身的一个子域，我们称之为平凡子域，其余与$F$不同的子域我们称为真子域。

定义：若域$F$中没有真子域，则称$F$为素域。

定理：给定域$F$，我们有：
\begin{enumerate}
\item 若$char(F)=p,p$为某一素数，则$F$包含的素域同构于$F_p$
\item 若$char(F)=0$，则$F$包含的素域同构于$Q$
\end{enumerate}

我们首先来证明第一条，考虑下面这个环同态：

\begin{align*}
\varphi:Z\mapsto F,\;\;\; n\mapsto n\cdot 1_F
\end{align*}

其中$1_F$指的是域$F$中的乘法单位元。由于$char(F)=p$，也就是说$p$是使$p\cdot 1_F=0$的最小正整数，因此：

\begin{align*}
ker\varphi=pZ
\end{align*}

那么由同态基本定理，就会有：

\begin{align*}
im\varphi \cong Z/pZ=F_p
\end{align*}

由于$p$是素数，那么$F_p$显然是一个域，由环同构得$im\varphi$也是一个域，所以其是$F$的一个子域，另外，$F$的任意子域肯定会包含乘法单位元$1_F$，那肯定会包含每个$n\cdot 1_F$，故包含$im\varphi$，所以$im\varphi$不会再有真子域了，因为$im\varphi$的子域肯定会是$F$的子域，但是$F$的所有子域都会包含$im\varphi$，所以$im\varphi$就是一个素域，同构于$F_p$

下面我们来证明第二条，我们依然考虑刚才的环同态$\varphi$，由于$char(F)=0$，所以只能有$ker\varphi=\{0\}$，所以对于任意的非零整数$b$，都会有$b\cdot 1_F\neq 0$，于是可以定义：

\begin{align*}
\psi:Q\mapsto F,\;\;\; \frac{a}{b}\mapsto (a\cdot 1_F)(b\cdot 1_F)^{-1}
\end{align*}

这个映射是良定义的，任意两个相同的原像都会映射到相同的像，另外，易得该映射能够保持加法乘法，并且：

\begin{align*}
\psi(\frac{a}{b})=0\Lightarrow a\cdot1_F=0\Lightarrow a=0
\end{align*}

由于只有0会被映射到0，所以当有$\psi(x)=\psi(y)$的时候，我们就能推出$\psi(x-y)=\psi(x)-psi(y)=0$，也就是会有$x=y$，因此$\psi$是单射，其象是$F$中一个同构于$Q$的子域。

最后，$F$的任意子域都包含整数倍的$a\cdot 1_F$，也包含非零元素$b\cdot 1_F$的逆元，因此包含$im \psi$，所以其就是$F$的素域，同构于$Q$。

定义：给定域$K<F$以及$M\subset F$，我们定义域$K(M)$为$F$中所有同时包含$K$和$M$的子域的交，并称之为由域$K$添加$M$中元素所得到的扩域

\begin{enumerate}
\item 若$M=\{\theta_1,\cdots,\theta_n\}$，我们也记$K(M)=K(\theta_1,\cdots,\theta_n)$。
\item 若$M=\{\theta\}$，则称$L=K(\theta)$为域$K$的单扩张，并称元素$\theta$为$L$在$K$上的一个定义元。
\item 可以验证$K(M)$是$F$中最小的同时包含$K$和$M$的子域。
\end{enumerate}

那这个定义有什么作用呢？其实就是想表示，在子域$K$中加入$M$的元素后，为了仍然是一个域，最少需要得到一个多大的域。

定义：给定域$K<F$以及$\theta\in F$，若$\theta$满足代数方程：

\begin{align*}
a_n\theta^n+a_{n-1}\theta^{n-1}+\cdots+a_1\theta+a_0=0
\end{align*}

其中$a_i\in K$并且不全为0，则称$\theta$在$K$上是代数的，如果不存在上面的代数方程，则称$\theta$在$K$上是超越的，如果$F$中的每一个元素在$K$上都是代数的，则称$F$是$K$的代数扩张，反之，若$F$中至少存在一个在$K$上超越的元素，则称$F$是$K$的超越扩张。

定理：假定元素$\theta\in F$是$K$上的代数元，令：

\begin{align*}
J=\{f\in K[x]:f(\theta)=0\}
\end{align*}

则存在唯一的首1多项式$g\in K[x]$使得$J=(g)$，并且多项式$g$在$K$上不可约。

如何证明这一定理呢？因为$\theta$是代数元，所以存在非零多项式$f$满足$f(\theta)=0$，所以$J$中有非零多项式，在这些多项式中，我们选一个次数最小的，再除以它的最高项系数，就能得到一个首1多项式$g$，并且仍然有$g(\theta)=0$，这是显然的，接下来我们证明$J=(g)$，任取$f\in J$，我们使用多项式带余除法：

\begin{align*}
f(x)=q(x)g(x)+r(x)
\end{align*}

其中$r(x)=0$或者$deg(r)<deg(g)$，带入$\theta$，就能得到：

\begin{align*}
0=f(\theta)=q(\theta)g(\theta)+r(\theta)
\end{align*}

这样的话就必然会有$r(\theta)=0$，也就是会有$r(x)\in J$，但是如果$r(x)$是非零多项式的话，其次数一定小于$g(x)$，这就违背了我们当时选取$g(x)$的要求，所以$r(x)$只能是0，所以$J$中任意元素都可以由$g$来生成，另一方面，$g$在$K[x]$上生成的任意元素显然都属于$J$，所以$J=(g)$。

如果$g$是可约的，不妨记为$g=uv$，那么：

\begin{align*}
0=g(\theta)=u(\theta)v(\theta)
\end{align*}    

由于$K$是一个域，所以当两元素乘积为0时，必定会有一个为0，也就是说$u$或者$v$会属于$J$，这也违背了我们当时选取$g(x)$的要求，所以说$g$是不可约的。最后证明唯一性：假设另一个首1多项式$h$也满足$J=(h)$，那么就会有：

\begin{align*}
g\in (h),h\in (g)
\end{align*}

那么就会有$g|h$并且$h|g$，所以这两个多项式次数必然是相等的，只相差一个常数倍，又因为都是首1多项式，所以常数就是1，也就是$g=h$，唯一性得证。

定义：我们称$g$为元素$\theta$在$K$上的极小多项式，也将$g$的次数称为元素$\theta$在$K$上的次数。

定理：假定元素$\theta\in F$是域$K$上的代数元，我们有以下关于$\theta$的极小多项式$g$的性质：

\begin{enumerate}
\item $g$是$K[x]$中的不可约多项式
\item 对于$f\in K[x]$，有$f(\theta)=0$当且仅当$g|f$
\item $g$是$K[x]$中以$\theta$为根的次数最低的首1多项式
\end{enumerate}

这三条性质在我们证明上一个定理时都证明过了，这里不再重复。

\section*{域扩张}


\end{document}